\documentclass[10pt,a4paper]{amsart}
\usepackage{amsmath,amssymb,amsthm,graphicx}
% Retained minimal preamble for the Stacks Project chapter more-algebra.tex (More on Algebra).
% The chapter file pinned-original-source.tex carries only \input{preamble} on its first line, so no
% preamble text lives inside the chapter itself.  The Stacks Project uses one single shared preamble for
% every chapter; the snapshot stages it identically as sources/stacks-chow/preamble.tex and
% sources/stacks-etale-cohomology/preamble.tex.  The block below is copied mechanically, line for line,
% from sources/stacks-chow/preamble.tex lines 14-18 (the Xy-pic setup) and 154-217 (the native theorem
% declarations together with the mathematical macros).  Nothing was added, removed or redefined; the
% wrapper class and the amsmath/amssymb/amsthm/graphicx base packages are supplied by the assembler.
% Chapter: more-algebra.tex -- More on Algebra.
% For commutative diagrams we use Xy-pic
\usepackage[all]{xy}

% We use 2cell for 2-commutative diagrams.
\xyoption{2cell}

% Theorem environments.
%
\theoremstyle{plain}
\newtheorem{theorem}[subsection]{Theorem}
\newtheorem{proposition}[subsection]{Proposition}
\newtheorem{lemma}[subsection]{Lemma}

\theoremstyle{definition}
\newtheorem{definition}[subsection]{Definition}
\newtheorem{example}[subsection]{Example}
\newtheorem{exercise}[subsection]{Exercise}
\newtheorem{situation}[subsection]{Situation}

\theoremstyle{remark}
\newtheorem{remark}[subsection]{Remark}
\newtheorem{remarks}[subsection]{Remarks}

\numberwithin{equation}{subsection}

% Macros
%
\def\lim{\mathop{\mathrm{lim}}\nolimits}
\def\colim{\mathop{\mathrm{colim}}\nolimits}
\def\Spec{\mathop{\mathrm{Spec}}}
\def\Hom{\mathop{\mathrm{Hom}}\nolimits}
\def\Ext{\mathop{\mathrm{Ext}}\nolimits}
\def\SheafHom{\mathop{\mathcal{H}\!\mathit{om}}\nolimits}
\def\SheafExt{\mathop{\mathcal{E}\!\mathit{xt}}\nolimits}
\def\Sch{\mathit{Sch}}
\def\Mor{\mathop{\mathrm{Mor}}\nolimits}
\def\Ob{\mathop{\mathrm{Ob}}\nolimits}
\def\Sh{\mathop{\mathit{Sh}}\nolimits}
\def\NL{\mathop{N\!L}\nolimits}
\def\CH{\mathop{\mathrm{CH}}\nolimits}
\def\proetale{{pro\text{-}\acute{e}tale}}
\def\etale{{\acute{e}tale}}
\def\QCoh{\mathit{QCoh}}
\def\Ker{\mathop{\mathrm{Ker}}}
\def\Im{\mathop{\mathrm{Im}}}
\def\Coker{\mathop{\mathrm{Coker}}}
\def\Coim{\mathop{\mathrm{Coim}}}

% Boxtimes
%
\DeclareMathSymbol{\boxtimes}{\mathbin}{AMSa}{"02}

%
% Macros for moduli stacks/spaces
%
\def\QCohstack{\mathcal{QC}\!\mathit{oh}}
\def\Cohstack{\mathcal{C}\!\mathit{oh}}
\def\Spacesstack{\mathcal{S}\!\mathit{paces}}
\def\Quotfunctor{\mathrm{Quot}}
\def\Hilbfunctor{\mathrm{Hilb}}
\def\Curvesstack{\mathcal{C}\!\mathit{urves}}
\def\Polarizedstack{\mathcal{P}\!\mathit{olarized}}
\def\Complexesstack{\mathcal{C}\!\mathit{omplexes}}
% \Pic is the operator that assigns to X its picard group, usage \Pic(X)
% \Picardstack_{X/B} denotes the Picard stack of X over B
% \Picardfunctor_{X/B} denotes the Picard functor of X over B
\def\Pic{\mathop{\mathrm{Pic}}\nolimits}
\def\Picardstack{\mathcal{P}\!\mathit{ic}}
\def\Picardfunctor{\mathrm{Pic}}
\def\Deformationcategory{\mathcal{D}\!\mathit{ef}}
\usepackage{cleveref}
\crefname{originallemma}{original lemma}{original lemmas}
\Crefname{originallemma}{Original lemma}{Original lemmas}
\makeatletter
\newlabel{algebra-lemma-Artin-Rees}{{00IN}{source}{}{sourcepointer-1}{}}
\newlabel{algebra-lemma-Artin-Rees@cref}{{[originallemma][0][]00IN}{[1][0][]source}{}{}{}}
\makeatother
\crefname{originallemma}{original lemma}{original lemmas}
\Crefname{originallemma}{Original lemma}{Original lemmas}
\makeatletter
\newlabel{algebra-lemma-map-AR}{{00IO}{source}{}{sourcepointer-2}{}}
\newlabel{algebra-lemma-map-AR@cref}{{[originallemma][0][]00IO}{[1][0][]source}{}{}{}}
\makeatother
\crefname{originallemma}{original lemma}{original lemmas}
\Crefname{originallemma}{Original lemma}{Original lemmas}
\makeatletter
\newlabel{categories-remark-pro-category}{{05PX}{source}{}{sourcepointer-3}{}}
\newlabel{categories-remark-pro-category@cref}{{[originallemma][0][]05PX}{[1][0][]source}{}{}{}}
\makeatother
\title{Artin-Rees controls the A-adic stabilisation of the homology of a complex of finite modules over a Noetherian ring}
\author{The Stacks Project Authors}
\date{stacks-project:more-algebra.tex:local-snapshot-20261012; GNU FDL 1.2 or later}
\begin{document}
\maketitle
\setcounter{subsection}{1}
\noindent\textit{Editorial note.} All delivered mathematical text is the author's own, copied from the pinned Stacks Project chapter file more-algebra.tex; no statement or proof step was rewritten, reordered or removed.\par\medskip
\noindent\textit{Editorial note.} The chapter file pinned-original-source.tex carries only \textbackslash{}input\{preamble\} on its first line, so no preamble text lives inside the chapter. The retained file sources/stacks-more-algebra/retained-preamble.tex is a mechanical, line-for-line copy of the project's single shared preamble as staged in the snapshot; nothing was added, removed or redefined and no package was added.\par\medskip
\noindent\textit{Editorial note.} The source URL is the source-level Stacks Project URL, as required when more than one card is produced from one source; the per-claim tag is carried by claim\_key.\par\medskip
\noindent\textit{Editorial note.} The author writes 'We omit the verification that the diagram commutes.' in the counted proof span; that sentence is kept verbatim and is the one step left to the reader. The author's typo 'the cass of z' in the last paragraph is reproduced unchanged.\par\medskip
\noindent\textit{Editorial note.} The three cross-chapter references inside the counted proof span (algebra-lemma-map-AR, algebra-lemma-Artin-Rees and categories-remark-pro-category) are recorded as non-load-bearing navigation pointers; both Artin-Rees applications and everything resting on them stay inside the counted proof span.\par\medskip
\noindent\textit{Allowed standard prerequisites (editorial).}\begin{itemize}
\item The Artin-Rees lemma for an ideal in a Noetherian ring. applied at proof lines 28916-28920 through the cited Algebra lemmas. Supplies the two constants c\_1 and c\_2 whose sum c controls every subsequent statement.
\item Pro-objects of modules and their morphisms. the isomorphism claim in part (1) at statement lines 28894-28895 and the cited Categories remark at proof lines 28962-28963. Fixes what it means for (H\_n) and (H/I\textasciicircum{}n H) to be isomorphic, which is the primary conclusion.
\item Inverse limits of modules and Mittag-Leffler conditions. parts (2) and (3) at statement lines 28895-28896. Makes sense of the limit comparison and of the stabilization of images claimed in part (4).
\end{itemize}\medskip
\ifdefined\hypertarget\hypertarget{sourcepointer-1}{}\fi
\noindent\textit{Source reference.} Original lemma 00IN: focused-seeds/AGP/sources/stacks-snapshot-20261011/tags.txt, source lines 685--685. Cross-chapter navigation pointer into the Stacks Project; the standard fact it supplies is applied inside the counted proof span and the argument resting on it stays there.\par\medskip
\ifdefined\hypertarget\hypertarget{sourcepointer-2}{}\fi
\noindent\textit{Source reference.} Original lemma 00IO: focused-seeds/AGP/sources/stacks-snapshot-20261011/tags.txt, source lines 686--686. Cross-chapter navigation pointer into the Stacks Project; the standard fact it supplies is applied inside the counted proof span and the argument resting on it stays there.\par\medskip
\ifdefined\hypertarget\hypertarget{sourcepointer-3}{}\fi
\noindent\textit{Source reference.} Original lemma 05PX: focused-seeds/AGP/sources/stacks-snapshot-20261011/tags.txt, source lines 7356--7356. Cross-chapter navigation pointer into the Stacks Project; the standard fact it supplies is applied inside the counted proof span and the argument resting on it stays there.\par\medskip
\begin{lemma}
\label{lemma-consequence-Artin-Rees}
Let $I$ be an ideal of a Noetherian ring $A$. Let
$
K \xrightarrow{\alpha} L \xrightarrow{\beta} M
$
be a complex of finite $A$-modules. Set $H = \Ker(\beta)/\Im(\alpha)$.
For $n \geq 0$ let
$$
K/I^nK \xrightarrow{\alpha_n} L/I^nL \xrightarrow{\beta_n} M/I^nM
$$
be the induced complex. Set $H_n = \Ker(\beta_n)/\Im(\alpha_n)$.
Then there are canonical $A$-module maps giving a commutative diagram
$$
\xymatrix{
& & & H \ar[lld] \ar[ld] \ar[d] \\
\ldots \ar[r] & H_3 \ar[r] & H_2 \ar[r] & H_1
}
$$
Moreover, there exists a $c > 0$ and canonical $A$-module maps
$H_n \to H/I^{n - c}H$ for $n \geq c$ such that the compositions
$$
H/I^n H \to H_n \to  H/I^{n - c}H
\quad\text{and}\quad
H_n \to H/I^{n - c}H \to H_{n - c}
$$
are the canonical ones. Moreover, we have
\begin{enumerate}
\item $(H_n)$ and $(H/I^nH)$ are isomorphic as pro-objects of $\text{Mod}_A$,
\item $\lim H_n = \lim H/I^n H$,
\item the inverse system $(H_n)$ is Mittag-Leffler,
\item the image of $H_{n + c} \to H_n$ is equal to the image of $H \to H_n$,
\item the composition $I^cH_n \to H_n \to H/I^{n - c}H \to H_n/I^{n - c}H_n$
is the inclusion $I^cH_n \to H_n$ followed by the quotient map
$H_n \to H_n/I^{n - c}H_n$, and
\item the kernel and cokernel of $H/I^nH \to H_n$ is annihilated by $I^c$.
\end{enumerate}
\end{lemma}

\begin{proof}
Observe that $H_n = \beta^{-1}(I^nM)/\Im(\alpha) + I^nL$. For $n \geq 2$
we have $\beta^{-1}(I^nM) \subset \beta^{-1}(I^{n - 1}M)$ and
$\Im(\alpha) + I^nL \subset \Im(\alpha) + I^{n - 1}L$. Thus
we obtain our canonical map $H_n \to H_{n - 1}$. Similarly, we have
$\Ker(\beta) \subset \beta^{-1}(I^nM)$ and
$\Im(\alpha) \subset \Im(\alpha) + I^nL$ which produces the
canonical map $H \to H_n$. We omit the verification that the diagram
commutes.

\medskip\noindent
By Artin-Rees we may choose $c_1, c_2 \geq 0$ such that
$\beta^{-1}(I^nM) \subset \Ker(\beta) + I^{n - c_1}L$ for $n \geq c_1$ and
$\Ker(\beta) \cap I^nL \subset I^{n - c_2}\Ker(\beta)$ for $n \geq c_2$, see
Algebra, Lemmas \ref{algebra-lemma-map-AR} and
\ref{algebra-lemma-Artin-Rees}. Set $c = c_1 + c_2$.

\medskip\noindent
Let $n \geq c$. We define $\psi_n : H_n \to H/I^{n - c}H$ as follows.
Say $x \in H_n$. Choose $y \in \beta^{-1}(I^nM)$ representing $x$.
Write $y = z + w$ with $z \in \Ker(\beta)$ and $w \in I^{n - c_1}L$
(this is possible by our choice of $c_1$). We set $\psi_n(x)$
equal to the class of $z$ in $H/I^{n - c}H$. To see this is well defined,
suppose we have a second set of choices $y', z', w'$ as above for $x$
with obvious notation.
Then $y' - y \in \Im(\alpha) + I^nL$, say $y' - y = \alpha(v) + u$
with $v \in K$ and $u \in I^nL$. Thus
$$
y' = z' + w' = \alpha(v) + u + z + w
\Rightarrow
z' = z + \alpha(v) + u + w - w'
$$
Since $\beta(z' - z - \alpha(v)) = 0$ we find that
$u + w - w' \in \Ker(\beta) \cap I^{n - c_1}L$
which is contained in $I^{n - c_1 - c_2}\Ker(\beta) = I^{n - c}\Ker(\beta)$
by our choice of $c_2$. Thus $z'$ and $z$ have the same image in
$H/I^{n - c}H$ as desired.

\medskip\noindent
The composition $H/I^n H \to H_n \to  H/I^{n - c}H$ is the canonical map
because if $z \in \Ker(\beta)$ represents an element $x$ in
$H/I^nH = \Ker(\beta)/\Im(\alpha) + I^n\Ker(\beta)$ then it is clear
from the above that $x$ maps to the class of $z$ in $H/I^{n - c}H$
under the maps constructed above.

\medskip\noindent
Let us consider the composition $H_n \to H/I^{n - c}H \to H_{n - c}$.
Given $x, y, z, w$ as in the construction of $\psi_n$ above, we see
that $x$ is mapped to the cass of $z$ in $H_{n - c}$. On the other hand,
the canonical map $H_n \to H_{n - c}$ from the first paragraph of the
proof sends $x$ to the class of $y$. Thus we have to show that
$y - z \in \Im(\alpha) + I^{n - c}L$ which is the case because
$y - z = w \in I^{n - c_1}L \subset I^{n - c}L$.

\medskip\noindent
Statements (1) -- (4) are formal consequences of what we just proved.
Namely, (1) follows from the existence of the maps and the definition
of morphisms of pro-objects in
Categories, Remark \ref{categories-remark-pro-category}.
Part (2) holds because isomorphic pro-objects have isomorphic limits.
Part (3) is immediate from part (4).
Part (4) follows from the factorization
$H_{n + c} \to H/I^nH \to H_n$ of the canonical map
$H_{n + c} \to H_n$.

\medskip\noindent
Proof of part (5). Let $x \in I^cH_n$. Write $x = \sum f_i x_i$ with
$x_i \in H_n$ and $f_i \in I^c$. Choose $y_i, z_i, w_i$ as in the
construction of $\psi_n$ for $x_i$. Then for the computation of
$\psi_n$ of $x$ we may choose
$y = \sum f_iy_i$, $z = \sum f_i z_i$ and $w = \sum f_i w_i$
and we see that $\psi_n(x)$ is given by the class of $z$.
The image of this in $H_n/I^{n - c}H_n$ is equal to the class
of $y$ as $w = \sum f_i w_i$ is in $I^nL$. This proves (5).

\medskip\noindent
Proof of part (6). Let $y \in \Ker(\beta)$ whose class is $x$ in $H$.
If $x$ maps to zero in $H_n$, then $y \in I^nL + \Im(\alpha)$.
Hence $y - \alpha(v) \in \Ker(\beta) \cap I^nL$ for some $v \in K$.
Then $y - \alpha(v) \in I^{n - c_2}\Ker(\beta)$ and hence the class
of $y$ in $H/I^nH$ is annihilated by $I^{c_2}$. Finally, let $x \in H_n$
be the class of $y \in \beta^{-1}(I^nM)$. Then we write $y = z + w$
with $z \in \Ker(\beta)$ and $w \in I^{n - c_1}L$ as above. Clearly, if
$f \in I^{c_1}$ then $fx$ is the class of
$fy + fw \equiv fy$ modulo $\Im(\alpha) + I^nL$ and hence $fx$ is
the image of the class of $fy$ in $H$ as desired.
\end{proof}

\end{document}
